\section{Cardinals}\label{sec:cardinals}

\paragraph{Equinumerosity and domination}

\begin{definition}\label{def:equinumerosity}\mcite[129]{Enderton1977SetTheory}
  We say that two sets are \term[bg=равномощни (множества) (\cite[79]{ГеновМиховскиМоллов1991Алгебра}), ru=равномощные (множества) (\cite[154]{КолмогоровДрагалин2006ЛогикаДополнительныеГлавы})]{equinumerous} if there exists a \hyperref[def:function_invertibility/bijective]{bijective function} between them.
\end{definition}
\begin{comments}
  \item As shown in \cref{thm:equinumerosity_is_equivalence}, equinumerosity can be regarded as an equivalence relation on the \hyperref[def:universal_set]{universal set} \( V \). The cardinal numbers, introduced in \cref{def:cardinal}, act as canonical representatives that are well-defined in \hyperref[def:zfc]{\logic{ZFC}}.

  \item Different terms for this relation can be found in the literature:
  \begin{itemize}
    \item \enquote{Equinumerous} is used in
    \cite[129]{Enderton1977SetTheory} and
    \cite[xxi]{Mendelson2015Logic}.

    \item \enquote{Equipotent} is used in
    \cite[def. 6.2.32(ii)]{Hinman2005Logic},
    \cite[def. III.3.1]{Bourbaki2004TheoryOfSets},
    \cite[3]{Harzheim2005OrderedSets} and
    \cite[3]{Engelking1989GeneralTopology}

    \item \enquote{Equipollent} is used in \cite[28]{Kelley1975GeneralTopology}.

    \item \enquote{Equivalent}, used in \cite[52]{Halmos1960NaiveSetTheory}.
  \end{itemize}
\end{comments}

\begin{proposition}\label{thm:equinumerosity_is_equivalence}
  \hyperref[def:equinumerosity]{Equinumerosity} satisfies the axioms for an \hyperref[def:equivalence_relation]{equivalence relation}.
\end{proposition}
\begin{proof}
  \SubProofOf*[def:binary_relation/reflexive]{reflexivity} Clearly the \hyperref[def:set_valued_function/identity]{identity} \( \id_A \) is a bijective endofunction on \( A \).

  \SubProofOf*[def:binary_relation/symmetric]{symmetry} If \( f: A \to B \) is bijective, so is its \hyperref[def:set_valued_function/inverse]{inverse} \( f^{-1}: B \to A \).

  \SubProofOf*[def:binary_relation/transitive]{transitivity} If \( f: A \to B \) and \( g: B \to C \) are bijective, \cref{thm:function_composition_invertibility} implies that their \hyperref[def:set_valued_function/composition]{composition} \( (g \bincirc f): A \to C \) is also bijective.
\end{proof}

\begin{definition}\label{def:set_domination}\mcite[145]{Enderton1977SetTheory}
  If there exists an \hyperref[def:function_invertibility/injective]{injective} function from \( A \) to \( B \) that is not necessarily \hyperref[def:function_invertibility/surjective]{surjective}, we say that \( A \) is \term{dominated by} \( B \) or that \( B \) \term{dominates} \( A \). If \( B \) dominates \( A \), and they are not equinumerous, we say that \( B \) \term[en=strictly dominates (\cite[90]{Halmos1960NaiveSetTheory})]{strictly dominates} \( A \).
\end{definition}
\begin{comments}
  \item We use set domination to define cardinal ordering in \cref{def:cardinal_ordering}. \Cref{thm:cardinal_ordering_is_well_ordering} shows that it is compatible with \hyperref[def:ordinal_ordering]{ordinal ordering}, i.e. set membership.

  \item As shown in \cref{thm:set_domination_is_preorder}, domination can be regarded as a preorder on the \hyperref[def:universal_set]{universal set} \( V \). \Fullref{thm:cantor_schroder_bernstein_theorem} shows that this order is compatible with equinumerosity.
\end{comments}

\begin{proposition}\label{thm:set_domination_is_preorder}
  \hyperref[def:set_domination]{Set domination} satisfies axioms of the theory of \hyperref[def:preordered_set]{preordered set}, i.e. \hyperref[def:binary_relation/reflexive]{reflexivity} and \hyperref[def:binary_relation/transitive]{transitivity}.
\end{proposition}
\begin{comments}
  \item Domination is not a partial order --- see \cref{ex:set_domination_and_equinumerosity/domination_without_equinumerosity}. It is \hyperref[def:binary_relation/connected]{connected}, however --- see \fullref{thm:set_domination_is_connected}.
\end{comments}
\begin{comments}
  \item We use this to prove that the cardinal ordering is a well-ordering in \cref{thm:cardinal_ordering_is_well_ordering}.
\end{comments}
\begin{proof}
  \SubProofOf*[def:binary_relation/reflexive]{reflexivity} Clearly the \hyperref[def:set_valued_function/identity]{identity} \( \id_A \) is an injective endofunction on \( A \), so \( A \) dominates itself.

  \SubProofOf*[def:binary_relation/transitive]{transitivity} If \( f: A \to B \) and \( g: B \to C \) are injective, \cref{thm:function_composition_invertibility} implies that their \hyperref[def:set_valued_function/composition]{composition} \( (g \bincirc f): A \to C \) is also injective.
\end{proof}

\begin{theorem}[Set domination is connected]\label{thm:set_domination_is_connected}
  For any two sets \( A \) and \( B \), either \( A \) \hyperref[def:set_domination]{dominates} \( B \) or vice versa.
\end{theorem}
\begin{comments}
  \item Within \hyperref[def:zfc]{\logic{ZF}}, this theorem is equivalent to the \hyperref[def:zfc/choice]{axiom of choice}. We list equivalent results in \cref{thm:axiom_of_choice_equivalences}.

  \item We use this to prove connectedness of cardinal ordering in \cref{thm:cardinal_ordering_is_well_ordering}.
\end{comments}
\begin{proof}
  We will show equivalence with \fullref{thm:well_ordering_theorem} rather than the axiom of choice.

  \ImplicationSubProof[thm:well_ordering_theorem]{well-ordering theorem}[thm:set_domination_is_connected]{domination trichotomy} Suppose that every set can be well-ordered.

  Let \( f: A \to \alpha \) and \( g: B \to \beta \) be order isomorphisms into their \hyperref[def:well_ordered_order_type]{order types}. \Cref{thm:ordinal_trichotomy} implies that \( \alpha \) and \( \beta \) are comparable, so we have three possibilities:
  \begin{itemize}
    \item If \( \alpha \in \beta \), then \( \alpha \subseteq \beta \) because \( \beta \) is transitive, so we have an inclusion function \( \iota: \alpha \to \beta \). \Cref{thm:function_composition_invertibility} implies that the composition \( g^{-1} \bincirc \iota \bincirc f: A \to B \) is injective.
    \item Similarly, if \( \beta \in \alpha \), then \( \beta \subseteq \alpha \) and thus \( f^{-1} \bincirc \iota \bincirc g: B \to A \) is injective.
    \item Otherwise, \( \alpha = \beta \) and \( f \bincirc g^{-1} \) and \( g \bincirc f^{-1} \) are mutual inverses.
  \end{itemize}

  Therefore, either \( A \) dominates \( B \) or vice versa.

  \ImplicationSubProof[thm:set_domination_is_connected]{domination trichotomy}[thm:well_ordering_theorem]{well-ordering theorem} Suppose that set domination is connected.

  Fix a set \( A \) and consider its \hyperref[def:hartogs_number]{Hartogs number} \( \alpha \). By \fullref{thm:hartogs_lemma}, there is no injective function from \( \alpha \) to \( A \), so there must be an injective function \( f: A \to \alpha \).

  It remains to define a well-ordering on \( A \) as
  \begin{equation*}
    x < y \T{if and only if} f(x) \in f(y).
  \end{equation*}
\end{proof}

\begin{example}\label{ex:set_domination_and_equinumerosity}
  We list examples related to \hyperref[def:set_domination]{set domination} and \hyperref[def:equinumerosity]{equinumerousity}:
  \begin{thmenum}
    \thmitem{ex:set_domination_and_equinumerosity/empty} \Cref{thm:empty_function_invertibility} implies that the empty function into any set is injective, so every set dominates the empty set.

    The same lemma implies that only the empty set with an empty codomain is surjective, so only the empty set is equinumerous with itself.

    \thmitem{ex:set_domination_and_equinumerosity/singleton} All \hyperref[def:singleton_set]{singleton sets} are equinumerous, and only the empty set is strictly dominated by singleton sets.

    \thmitem{ex:set_domination_and_equinumerosity/domination_without_equinumerosity} According to \cref{thm:set_domination_is_preorder}, set domination is a preorder. It is not a partial order, however, because antisymmetry can fail.

    For example, the map \( n \mapsto 2n \) from all natural numbers \( \BbbN \) to the even natural numbers \( 2\BbbN \) is injective and the identity map on \( 2\BbbN \) is an injective function from \( 2\BbbN \) to \( \BbbN \), however \( \BbbN \) and \( 2\BbbN \) are not equal.

    They are, however, equinumerous, because we can send every even number \( 2n \) to \( n \).

    This conclusion is generalized in \fullref{thm:cantor_schroder_bernstein_theorem}.

    \thmitem{ex:set_domination_and_equinumerosity/omega_plus_one} The \hyperref[def:smallest_inductive_set]{smallest inductive set} \( \omega \) and its \hyperref[def:ordinal_successor]{successor} \( S(\omega) \) are equinumerous, as demonstrated by the map
    \begin{equation*}
      \begin{aligned}
        &f: \omega \to S(\omega), \\
        &f(n) \coloneqq \begin{cases}
          \omega, &n = \varnothing, \\
          m,      &n = S(m).
        \end{cases}
      \end{aligned}
    \end{equation*}

    This example is important for the definition of initial ordinals in \cref{def:initial_ordinal} and of cardinals in \cref{def:cardinal}.
  \end{thmenum}
\end{example}

\begin{lemma}\label{thm:three_equinumerous_sets_lemma}\mcite[156]{КолмогоровДрагалин2006ЛогикаДополнительныеГлавы}
  If \( A \subseteq B \subseteq C \) and \( A \) is \hyperref[def:equinumerosity]{equinumerous} with \( C \), then \( B \) is equinumerous with \( C \).
\end{lemma}
\begin{proof}
  If \( B = C \), the lemma is trivial since the \hyperref[def:set_valued_function/identity]{identity function} \( \id_B: B \to B \) is bijective. If \( B \subsetneq C \), however, the identity \( \id_B \) must be extended in order to be a bijective function between \( B \) and \( C \). It will actually be simpler for us to define a function from \( C \) to \( B \).

  Let \( f: C \to A \) be a bijective function (such a function exists by the statement of the lemma). Consider the collection
  \begin{equation*}
    I \coloneqq \bigcap\set{ X \subseteq C \given (C \setminus B) \subseteq X \T{and} f[X] \subseteq X }
  \end{equation*}
  of all intermediate sets between \( C \setminus B \) and \( C \) that are invariant under \( f \). Consider the function
  \begin{equation*}
    \begin{aligned}
      &g: C \to B \\
      &g(x) \coloneqq \begin{cases}
        x,    &x \in C \setminus I, \\
        f(x), &x \in I.
      \end{cases}
    \end{aligned}
  \end{equation*}

  \SubProof{Proof of well-definedness} We must show that the image of \( g \) is contained in \( B \), as specified. Indeed, \( g \) is well-defined because, by \fullref{thm:de_morgans_laws_for_sets}, we have
  \begin{align*}
    C \setminus I
    &=
    C \setminus \bigcap\set{ X \subseteq C \given (C \setminus B) \subseteq X \T{and} f[X] \subseteq X }
    = \\ &=
    \bigcup\set{ C \setminus X \given X \subseteq C \T{and} (C \setminus B) \subseteq X \T{and} f[X] \subseteq X }.
  \end{align*}

  If \( (C \setminus B) \subseteq X \subseteq C \), \cref{thm:set_difference/complement} implies that
  \begin{equation*}
    C \setminus X \subseteq \underbrace{C \setminus (C \setminus B)}_B,
  \end{equation*}
  so the entire union above is a subset of \( B \).

  \SubProof{Proof of invariance} Both \( I \) and \( C \setminus I \) are invariant under \( g \). Indeed, \( C \setminus I \) is invariant because the restriction of \( g \) is the identity. On the other hand, if \( x \) is in \( I \), then \( x \) belongs to \( X \) for every subset \( X \) of \( C \) containing \( C \setminus B \) and invariant under \( f \). Thus, \( f(x) \) belongs to every such subset, and thus to their intersection, \( I \).

  \SubProofOf[def:function_invertibility/injective]{injectivity} Let \( g(x_1) = g(x_2) \) for some members \( x_1 \) and \( x_2 \) of \( C \). As already shown, both sets \( I \) and \( C \setminus I \) are invariant under \( f \), so both values \( x_1 \) and \( x_2 \) belong to either \( I \) or \( C \setminus I \). This leaves us with two possibilities:
  \begin{itemize}
    \item If \( x_1 \) and \( x_2 \) both belong to \( I \), then \( f(x_1) = g(x_1) = g(x_2) = f(x_2) \), hence \( x_1 = x_2 \) because \( f \) is bijective.

    \item If \( x_1 \) and \( x_2 \) both belong to \( C \setminus I \), then \( x_1 = g(x_1) = g(x_2) = x_2 \).
  \end{itemize}

  \SubProofOf[def:function_invertibility/surjective]{surjectivity} Suppose that \( g \) is not surjective and fix a value \( y \) from \(B \setminus g[C] \).

  \begin{itemize}
    \item If \( y \) is in \( I \), by invariance \( g(y) \) also is. This contradicts the definition of \( g \) since \( B \) is disjoint from \( I \).

    \item Otherwise, by invariance \( g(y) \) belongs to \( C \setminus I \), so \( g(y) = y \) belongs to \( g[C] \). This contradicts our choice of \( y \).
  \end{itemize}

  The obtained contradictions show that \( g \) is surjective.
\end{proof}

\begin{theorem}[Cantor-Schr\"oder-Bernstein theorem]\label{thm:cantor_schroder_bernstein_theorem}
  If two sets \hyperref[def:set_domination]{dominate} each other, they are \hyperref[def:equinumerosity]{equinumerous}.
\end{theorem}
\begin{comments}
  \item Our proof is based on \cite[156]{КолмогоровДрагалин2006ЛогикаДополнительныеГлавы} and is mostly delegated to \cref{thm:three_equinumerous_sets_lemma}.

  \item \Textcite[147]{Enderton1977SetTheory} explains the (early) history of this theorem. Generally, different authors use variations of its name:
  \begin{itemize}
    \item Cantor, Schr\"oder and Bernstein are all attributed in
    \cite[3]{Harzheim2005OrderedSets},
    \cite[thm. 2.22]{Гуров2013ТеорияРешёток},
    \cite[156]{КолмогоровДрагалин2006ЛогикаДополнительныеГлавы} and
    \cite[thm. 8.7.3]{Тыртышников2017ОсновыАлгебры}.

    \item Only Schr\"oder and Bernstein are attributed in
    \cite[thm. I.10.2]{Kunen2013SetTheory},
    \cite[thm. 6.4.4]{Hinman2005Logic},
    \cite[147]{Enderton1977SetTheory} and
    \cite[thm. 2.5.3]{Rosen2019DiscreteMathematics}.

    \item Only Cantor and Bernstein are attributed in
    \cite[thm. 3.2]{Jech2006SetTheory},
    \cite[4]{Engelking1989GeneralTopology} and
    \cite[15]{Willard2004GeneralTopology}.

    \item Only Bernstein is attributed in
    \cite[prop. 4.23(d)]{Mendelson2015Logic}.
  \end{itemize}
\end{comments}
\begin{proof}
  Let \( f: A \to B \) and \( g: B \to A \) be injective functions. From \cref{thm:function_composition_invertibility} it follows that \( g \bincirc f: A \to A \) is also an injective function. If we restrict its codomain to its image \( g[f[A]] \), it becomes bijective. Hence, \( A \) is equinumerous with \( g[f[A]] \). Since \( g[f[A]] \subseteq g[B] \subseteq A \), from \cref{thm:three_equinumerous_sets_lemma} it follows that \( A \) is equinumerous with \( g[B] \), which is the desired result.
\end{proof}

\begin{theorem}[Cantor's power set theorem]\label{thm:cantors_power_set_theorem}\mcite[thm. 6.4.10]{Hinman2005Logic}
  The \hyperref[def:power_set]{power set} of any set \( A \) \hyperref[def:set_domination]{strictly dominates} \( A \).
\end{theorem}
\begin{proof}
  The map \( x \mapsto \set{ x } \) from \( A \) to \( \pow(A) \) is clearly injective, therefore \( \pow(A) \) dominates \( A \). The converse is not true, however.

  Indeed, suppose that \( A \) dominates \( \pow(A) \). Then \fullref{thm:cantor_schroder_bernstein_theorem} implies that they are equinumerous, so so there exists some bijective function \( f: A \to \pow(A) \). Let
  \begin{equation*}
    B \coloneqq \set{ x \in A \colon x \notin f(x) }.
  \end{equation*}

  As a subset of \( A \), \( B \) belongs to \( \pow(A) \), and hence has a preimage under \( f \), say \( b \). If \( b \in B \), then \( b \notin f(b) \) by the definition of \( B \). Conversely, If \( b \notin f(b) \), then \( b \in B \).

  Thus, \( b \in B \) if and only if \( b \notin B \) --- a contradiction. Therefore, \( A \) cannot dominate \( \pow(A) \).
\end{proof}

\begin{theorem}[Cantor's paradox]\label{thm:cantors_paradox}\mcite[101]{Halmos1960NaiveSetTheory}
  In both \hyperref[def:naive_set_theory]{naive set theory} and \hyperref[def:zfc]{\logic{ZFC}}, the existence of the \hyperref[def:universal_set]{universal set} leads to a contradiction.
\end{theorem}
\begin{comments}
  \item Since the universal set can be defined using \hyperref[def:naive_set_theory/comprehension]{unrestricted comprehension}, this shows that naive set theory is \hyperref[def:fol_theory/consistent]{inconsistent}.

  \item This paradox arises from self-reference; see \cref{ex:self_reference/russells_paradox}. It is related to \fullref{thm:russels_paradox}, however we state it less formally.
\end{comments}
\begin{proof}
  Suppose that the universal set \( V \) exists. \Fullref{thm:cantors_power_set_theorem} implies that \( \pow(V) \) strictly dominates \( V \), so no injective functions from \( \pow(V) \) to \( V \) exist.

  Since \( V \) contains all sets, however, \( \pow(V) \) must be a subset of \( V \), so we have an injective inclusion function \( \iota: \pow(V) \to A \).

  We have obtained a contradiction, as desired.
\end{proof}

\paragraph{Initial ordinals}

\begin{definition}\label{def:initial_ordinal}\mcite[199]{Enderton1977SetTheory}
  If no member of the \hyperref[def:ordinal]{ordinal} \( \alpha \) is \hyperref[def:equinumerosity]{equinumerous} with it, we call \( \alpha \) an \term[ru=initial ordinal (\cite[199]{Enderton1977SetTheory})]{initial ordinal}.
\end{definition}
\begin{comments}
  \item Extensionally, ordinals are formally the same as cardinals (which we define in \cref{def:cardinal}). Intentionally, however, they are different --- see \cref{rem:ordinals_and_cardinals}.

  \item \Cref{thm:cardinal_is_infinite_iff_limit_ordinal} implies that infinite successor ordinals are not initial. A more explicit example is that \( S(\omega) \) is not initial --- see \cref{ex:set_domination_and_equinumerosity/omega_plus_one}
\end{comments}

\begin{proposition}\label{thm:def:initial_ordinal}
  \hyperref[def:initial_ordinal]{Initial ordinals} have the following basic properties:
  \begin{thmenum}
    \thmitem{thm:def:initial_ordinal/from_general} If a set is \hyperref[def:equinumerosity]{equinumerous} with an ordinal, it is equinumerous with an initial ordinal.

    \thmitem{thm:def:initial_ordinal/union} The union of initial ordinals is again an initial ordinal.
  \end{thmenum}
\end{proposition}
\begin{proof}
  \SubProofOf{thm:def:initial_ordinal/from_general} Suppose that \( A \) is equinumerous with the ordinal \( \alpha \). If \( \alpha \) is not itself initial, then \( A \) is equinumerous with some member of \( \alpha \). Since \( \alpha \) is well-ordered, we can take the smallest such member, which would be an initial ordinal.

  \SubProofOf{thm:def:initial_ordinal/union} \Cref{thm:def:ordinal_ordering/supremum} implies that the union \( \bigcup A \) of a set \( A \) of initial ordinals is their supremum with respect to set membership.

  To show that it is initial, suppose that there exists a bijective function \( f: \bigcup A \to \beta \) for some member \( \beta \) of \( \bigcup A \).

  Then \( \beta \) is a member of some \( \alpha \) in \( A \), hence it is a subset of \( \alpha \), and we have an inclusion function \( \iota: \beta \to \alpha \). The composition \( \iota \bincirc f: \bigcup A \to \alpha \) is injective by \cref{thm:function_composition_invertibility}. Furthermore, since \( \alpha \) is a member of \( A \), \cref{thm:unary_set_union_inequality} implies that \( \alpha \) is a subset of \( \bigcup A \). \Fullref{thm:cantor_schroder_bernstein_theorem} implies that \( \alpha \) and \( \bigcup A \) are equinumerous, and hence so are \( \alpha \) and \( \beta \). But \( \alpha \) is initial, so it cannot be equinumerous with a member of itself.

  The obtained contradiction implies that \( \bigcup A \) must also be initial ordinal.
\end{proof}

\begin{proposition}\label{thm:initial_ordinal_existence}
  Assuming the \hyperref[def:zfc/choice]{axiom of choice}, every set is \hyperref[def:equinumerosity]{equinumerous} with a unique \hyperref[def:initial_ordinal]{initial ordinal}.
\end{proposition}
\begin{proof}
  \ExistenceSubProof The axiom of choice implies \fullref{thm:well_ordering_theorem}, which gives us a well-ordering \( {<} \) for every set \( A \). The \hyperref[def:well_ordered_order_type]{order type} \( \ord(A, <) \) is then an ordinal equinumerous with \( A \). Based on \( \ord(A, <) \), \cref{thm:def:initial_ordinal/from_general} then gives us an initial ordinal equinumerous with \( A \).

  \UniquenessSubProof Let \( \alpha \) and \( \beta \) be initial ordinals. If \( f: A \to \alpha \) and \( g: A \to \beta \) are bijective functions, then, by \cref{thm:function_composition_invertibility}, so is their composition \( g \bincirc f^{-1}: \alpha \to \beta \).

  By \cref{thm:ordinal_trichotomy}, \( \alpha \) and \( \beta \) are comparable, which gives us three options:
  \begin{itemize}
    \item If \( \beta \in \alpha \), then \( \alpha \) cannot be initial.
    \item If \( \alpha \in \beta \), then \( \beta \) cannot be initial.
    \item It remains for them to be equal.
  \end{itemize}
\end{proof}

\begin{lemma}\label{thm:natural_number_is_not_equinumerous_to_proper_subset}
  No natural number (as a member of \hyperref[def:smallest_inductive_set]{\( \omega \)}) is equinumerous to a proper subset of itself.
\end{lemma}
\begin{proof}
  We will use \fullref{thm:omega_induction} on \( n \in \omega \):
  \begin{itemize}
    \item The case \( n = \varnothing \) is vacuous.
    \item Suppose that \( n \) is not equinumerous with a proper subset of itself. Fix any injective endofunction \( f: S(n) \to S(n) \). We must show that it is surjective.

    Suppose that it is not, and let \( m \) be the smallest set in \( S(n) \setminus f[S(n)] \).
    \begin{itemize}
      \item If \( m = n \), then \( n \) is an \hyperref[def:invariant_subset]{invariant subset} under \( f \), which allows us to consider the restriction \( f\restr_n: n \to n \). \Cref{thm:restriction_of_injective_functions} implies that \( f\restr_n \) is injective. It is also surjective since otherwise we would contradict the inductive hypothesis.

      But then \( f \) must map \( n \) to some value in \( n \), which it cannot because it is injective and \( n \) is already in the image of \( f\restr_n \). This contradicts the existence of \( m \).

      \item Otherwise, \( m \in n \), so we consider the function
      \begin{equation*}
        \begin{aligned}
          &g: n \to n, \\
          &g(k) \coloneqq \begin{cases}
            m,    &f(k) = n, \\
            f(k), &\T{otherwise}.
          \end{cases}
        \end{aligned}
      \end{equation*}

      This function is injective because \( g(k_1) = g(k_2) \) leads to \( f(k_1) = f(k_2) \) in all possible cases. Again by the inductive hypothesis, it is surjective, so \( f \) cannot be injective.
    \end{itemize}

    It both cases, we obtain a contradiction, so \( S(n) \setminus f[S(n)] \) must be empty. Therefore, \( f \) is surjective, and hence bijective.
  \end{itemize}
\end{proof}

\begin{proposition}\label{thm:members_of_omega_are_initial_ordinals}
  The natural numbers, as members of \hyperref[def:smallest_inductive_set]{smallest inductive set} \( \omega \), are \hyperref[def:initial_ordinal]{initial ordinals}.
\end{proposition}
\begin{comments}
  \item We restate this for \hyperref[def:cardinal]{cardinals} in \cref{thm:members_of_omega_are_initial_ordinals}.
\end{comments}
\begin{proof}
  Fix some \( n \in \omega \). \Cref{thm:ordinal_subset_inclusion} implies that every member of \( n \) is a also proper subset of \( n \). From \cref{thm:natural_number_is_not_equinumerous_to_proper_subset} it follows that no ordinal strictly smaller than \( n \) is equinumerous with \( n \).
\end{proof}

\begin{definition}\label{def:omega_hierarchy}\mcite[30]{Jech2006SetTheory}
  We use the unbounded transfinite recursion trick from \cref{rem:unbounded_transfinite_recursion} to define an unbounded family of infinite ordinals indexed by \( \alpha \):
  \begin{equation*}
    \omega_\alpha \coloneqq \begin{cases}
      \omega_0,                                                           &\alpha = \varnothing, \\
      \T{\hyperref[def:hartogs_number]{Hartogs' number} of} \omega_\beta, &\alpha = S(\beta), \\
      \bigcup\set{ \omega_\beta \given \beta \in \alpha }                 &\alpha \T{is a nonzero limit ordinal.}
    \end{cases}
  \end{equation*}
\end{definition}
\begin{comments}
  \item We use this hierarchy to define the \( \aleph \) hierarchy in \cref{def:aleph_hierarchy}.
\end{comments}

\begin{proposition}\label{thm:def:omega_hierarchy}
  The \( \omega \) hierarchy from \cref{def:omega_hierarchy} has the following basic properties:
  \begin{thmenum}
    \thmitem{thm:def:omega_hierarchy/monotone} It is monotone: \( \beta \in \alpha \) implies that \( \omega_\beta \in \omega_\alpha \).
  \end{thmenum}
\end{proposition}
\begin{proof}
  \SubProofOf{thm:def:omega_hierarchy/monotone} We will use \fullref{thm:transfinite_induction} on \( \alpha \).
  \begin{itemize}
    \item The base case \( \alpha = \varnothing \) is vacuous.
    \item Suppose that \( \alpha = S(\gamma) \) and that the inductive hypothesis holds for \( \gamma \). Fix some \( \beta \in \alpha \).

    In this case \( \omega_\alpha \) is the Hartogs' number of \( \omega_\gamma \), so \( \omega_\gamma \) belongs to \( \omega_\alpha \) by definition. This handles the case \( \beta = \gamma \).

    If \( \beta \in \gamma \), then \( \omega_\beta \) belongs to \( \omega_\gamma \) by the inductive hypothesis, and hence to \( \omega_\alpha \) because the latter is transitive.

    \item Finally, suppose that \( \alpha \) is a nonzero limit ordinal and that the inductive hypothesis holds for all members of \( \alpha \).

    In this case \( \omega_\alpha = \bigcup A \), where
    \begin{equation*}
      A \coloneqq \set{ \omega_\gamma \given \gamma \in \alpha }.
    \end{equation*}

    \Cref{thm:def:ordinal_ordering/supremum} implies that \( \omega_\alpha \) is the supremum of \( A \). Furthermore, \( \omega_\alpha \) does not itself belong to \( A \) because \( \alpha \) is a limit ordinal and \( \omega_\alpha \in A \) would imply \( \omega_{S(\alpha)} \in A \).

    Therefore, \( \omega_\beta \in \omega_\alpha \) for every \( \beta \in \alpha \).
  \end{itemize}
\end{proof}

\begin{lemma}\label{thm:omega_is_supremum_of_finite_ordinals}
  The \hyperref[def:smallest_inductive_set]{smallest inductive set} \( \omega \) is the union of its members, i.e. \( \omega = \bigcup \omega \).
\end{lemma}
\begin{proof}
  For every \( n \in \omega \), Since \( S(n) \in \omega \), \cref{thm:unary_set_union_inequality} implies that \( S(n) \subseteq \bigcup \omega \), hence \( n \in \bigcup \omega \).

  Conversely, for every \( n \in \bigcup \omega \), there exists some \( m \in \omega \) such that \( n \in m \). But \( m \subseteq \omega \) since \( \omega \) is transitive, so \( n \in \omega \).
\end{proof}

\begin{lemma}\label{thm:hartogs_number_is_initial_ordinal}
  \hyperref[def:hartogs_number]{Hartogs' numbers} are \hyperref[def:initial_ordinal]{initial ordinals}.
\end{lemma}
\begin{proof}
  Let \( \alpha \) be the Hartogs' number of the set \( A \).

  Suppose that there exists a bijective function \( f: \alpha \to \beta \) for some ordinal \( \beta \in \alpha \). Then \cref{thm:hartogs_lemma/minimal} implies that there is also an injective function \( g: \beta \to A \), so composing them gives us an injective function \( g \bincirc f: \alpha \to A \). But this contradicts \cref{thm:hartogs_lemma/injective}.

  Therefore, \( \alpha \) is not equinumerous with any smaller ordinal, i.e. it is an initial ordinal.
\end{proof}

\begin{proposition}\label{thm:omega_hierarchy_ordinals_iff_initial}
  An ordinal belongs to the \( \omega \) hierarchy from \cref{def:omega_hierarchy} if and only if it is an \hyperref[def:initial_ordinal]{initial ordinal} containing \( \omega_0 \) as a subset.
\end{proposition}
\begin{proof}
  \SufficiencySubProof We will use \fullref{thm:transfinite_induction} on \( \alpha \) to show that \( \omega_\alpha \) is an initial ordinal.
  \begin{itemize}
    \item For the base case \( \omega = \omega_0 \), \cref{thm:omega_is_supremum_of_finite_ordinals} implies \( \omega = \bigcup \omega \). \Cref{thm:members_of_omega_are_initial_ordinals} implies that the members of \( \omega \) are initial ordinals, so, by \cref{thm:def:initial_ordinal/union}, their union \( \omega \) is also initial.

    \item Irrespective of whether \( \omega_\alpha \) is initial, \cref{thm:hartogs_number_is_initial_ordinal} implies that, as a Hartogs' number, \( \omega_{S(\alpha)} \) is initial.

    \item Finally, if \( \set{ \omega_\beta \given \beta \in \alpha } \) are initial ordinals, by \cref{thm:def:initial_ordinal/union}, so is their union.
  \end{itemize}

  \NecessitySubProof We will use \fullref{thm:transfinite_induction} on \( \alpha \) to show that if \( \alpha \) is an initial ordinal containing \( \omega_0 \), then \( \alpha = \omega_\beta \) for some \( \beta \).

  Fix some initial ordinal \( \alpha \) and let
  \begin{equation*}
    \beta \coloneqq \set{ \gamma \given \omega_\gamma \in \alpha }.
  \end{equation*}

  This is a set by the \hyperref[def:zfc/replacement]{axiom (schema) of replacement}. Furthermore, it is a transitive set because, by \cref{thm:def:omega_hierarchy/monotone}, \( \delta \in \gamma \) implies \( \omega_\delta \in \omega_\gamma \), thus \( \omega_\delta \in \alpha \) and \( \delta \in \beta \).

  Then \cref{thm:def:ordinal/transitive_set} implies that \( \beta \) is itself an ordinal. We claim that \( \alpha = \omega_\beta \).
  \begin{itemize}
    \item If \( \beta = \varnothing \), then no members of the \( \omega \) sequence belong to \( \alpha \). Since \( \omega_0 \) is a subset of \( \alpha \), if \( \alpha \) is strictly larger, \cref{thm:ordinal_subset_inclusion} implies that \( \omega_0 \in \alpha \), which is a contradiction.

    It follows that \( \beta = \varnothing \) implies that \( \alpha = \omega_0 = \omega_\beta \).

    \item If \( \beta = S(\gamma) \), then \( \gamma \in \beta \) and thus \( \omega_\gamma \in \alpha \) by the definition of \( \beta \). But \( \omega_\beta \) cannot belong to \( \alpha \) since that would imply \( \beta \in \beta \), so \( \alpha \) is the smallest initial ordinal not containing \( \omega_\beta \). Since \( \omega_\beta \) is by assumption initial, it follows that \( \alpha = \omega_\beta \).

    \item If \( \beta \) is a nonzero limit ordinal, then \( \omega_\beta = \bigcup A \), where
    \begin{equation*}
      A = \set{ \omega_\gamma \given \gamma \in \beta }.
    \end{equation*}

    But \( \gamma \in \beta \) if and only if \( \omega_\gamma \in \alpha \), so
    \begin{equation*}
      A = \set{ \omega_\gamma \given \omega_\gamma \in \alpha }.
    \end{equation*}

    \Cref{thm:def:ordinal_ordering/supremum} implies that \( \omega_\beta \) is the supremum of all \( \omega_\gamma \) in \( \alpha \), and \cref{thm:def:initial_ordinal/union} implies that \( \omega_\beta \) is initial. Furthermore, \( \omega_\beta \) cannot belong to \( \alpha \) since that would lead to \( \beta \in \beta \).

    Thus, as in the previous case, \( \omega_\beta \) is the smallest initial ordinal not in \( \alpha \), so \( \alpha = \omega_\beta \).
  \end{itemize}
\end{proof}

\paragraph{Cardinality}

\begin{remark}\label{rem:ordinals_and_cardinals}
  As shown in \cref{thm:equinumerosity_is_equivalence}, \hyperref[def:equinumerosity]{equinumerousity} acts as an \hyperref[def:equivalence_relation]{equivalence relation} on sets. We want to introduce a concept of cardinality that characterizes equinumerosity --- an operator \( \card \) on sets that satisfies \fullref{thm:humes_principle}:
  \begin{equation*}
    \card(A) = \card(B) \T{if and only if} A \T{and} B \T{are equinumerous.}
  \end{equation*}

  This requires the value \( \card(A) \) to be a set. A naive attempt would be to use the equivalence class itself as a value for \( \card(A) \), however all equivalence classes except that for \( \varnothing \) are \hyperref[con:set_theory_class]{proper classes}.

  \Cref{thm:initial_ordinal_existence} implies that, assuming the \hyperref[def:zfc/choice]{axiom of choice}, every set \( A \) is equinumerous with a unique \hyperref[def:initial_ordinal]{initial ordinal} \( \alpha \). This prompts the common definition of cardinality via initial ordinals, which we give in \cref{def:cardinality}. We define cardinal numbers in \cref{def:cardinal} as initial ordinals since they are precisely the possible cardinalities.

  It is important, however, to distinguish ordinals and cardinals. They have different notions of \enquote{successors}, as highlighted in \cref{thm:cardinal_is_infinite_iff_limit_ordinal} --- all infinite successor cardinals are limit ordinals. We also define arithmetic operations differently on them --- see \fullref{sec:ordinal_arithmetic} and \fullref{sec:cardinal_arithmetic}. Therefore, initial ordinals and cardinals are identical extensionally, but not intentionally. In \cite[ch. 10]{UnivalentFoundationsProgram2013HoTT}, this is made concrete by introducing different \hyperref[def:type_assertion]{types} for ordinals and cardinals.
\end{remark}
\begin{comments}
  \item An alternative approach to cardinality is to use a trick due to Dana Scott to define cardinals using the \hyperref[def:zfc/foundation]{axiom of foundation} via \fullref{thm:axiom_of_regularity} instead of the axiom of choice via \fullref{thm:well_ordering_theorem}. It does not lead to a satisfactory notions of cardinals, however. This is discussed in \cref{rem:cardinals_via_scotts_trick}.
\end{comments}

\begin{definition}\label{def:cardinality}\mimprovised
  As shown in \cref{thm:initial_ordinal_existence}, assuming the \hyperref[def:zfc/choice]{axiom of choice}, every set \( A \) is \hyperref[def:equinumerosity]{equinumerous} with a unique \hyperref[def:initial_ordinal]{initial ordinal} \( \alpha \). We call \( \alpha \) the \term[en=cardinality (\cite[147]{Enderton1977SetTheory}), ru=мощность (множества) (\cite[\S 1.5.13]{КусраевКутателадзе2005БулевозначныйАнализ})]{cardinality} of \( A \) and denote it by \( \card(A) \).
\end{definition}
\begin{comments}
  \item Cardinality serves as a quantification of equinumerosity, and the particular definition is not as important as the fact that it satisfies \fullref{thm:humes_principle}. We discuss this more broadly in \cref{rem:ordinals_and_cardinals}.
\end{comments}

\begin{theorem}[Hume's principle]\label{thm:humes_principle}\mcite{Shapiro2008MathematicalObjects}
  Two sets are \hyperref[def:equinumerosity]{equinumerous} if and only if they have the same \hyperref[def:cardinality]{cardinality}.
\end{theorem}
\begin{comments}
  \item This principle, attributed to David Hume by Gottlob Frege, predates modern formal logic. It is stated in \cite[168]{Shapiro2008MathematicalObjects} as follows:
  \begin{displayquote}
    For any concepts \( F \), \( G \), the number of \( F \)'s is identical to the number of \( G \)'s if and only if \( F \) is equinumerous with \( G \).
  \end{displayquote}
\end{comments}
\begin{proof}
  We have shown in \Cref{thm:equinumerosity_is_equivalence} that equinumerosity is transitive, Therefore, \( A \) and \( B \) are equinumerous if and only if they are both equinumerous with the same cardinal.
\end{proof}

\begin{definition}\label{def:cardinal}\mcite[29]{Jech2006SetTheory}
  We refer to the possible \hyperref[def:cardinality]{cardinalities} as \term{cardinal numbers} or simply \term[ru=кардинал (\cite[\S 1.5.12]{КусраевКутателадзе2005БулевозначныйАнализ})]{cardinals}.

  Extensionally, these are simply the \hyperref[def:initial_ordinal]{initial ordinals}, however intentionally ordinals and cardinals must be distinguished because they have conflicting operations --- see \cref{rem:ordinals_and_cardinals}. We try not to rely on the encoding of cardinals as initial ordinals.
\end{definition}
\begin{comments}
  \item Within the metalanguage, we denote abstract ordinals via the Greek letters \( \kappa \), \( \lambda \), \( \mu \) and \( \nu \), in accordance with \cref{rem:mathematical_logic_conventions/greek_alphabet}.

  \item We discuss in \cref{ex:def:category_skeleton/cardinals} how the \hyperref[def:small_set]{small} \hyperref[def:cardinal]{cardinal numbers} form a \hyperref[def:category_skeleton]{skeleton} of the \hyperref[def:category_of_small_sets]{category of small sets}.
\end{comments}

\paragraph{Cardinal ordering}

\begin{definition}\label{def:cardinal_ordering}\mcite[145]{Enderton1977SetTheory}
  For \hyperref[def:cardinal]{cardinals}, we introduce the following order relation:
  \begin{equation*}
    \kappa \leq \mu \T{if and only if} \kappa \T{is \hyperref[def:set_domination]{dominated} by} \mu.
  \end{equation*}
\end{definition}

\begin{proposition}\label{thm:cardinal_ordering_is_well_ordering}
  \hyperref[def:cardinal_ordering]{Cardinal ordering} is a \hyperref[def:well_ordered_set]{well-ordering}, i.e. it satisfies the axioms for well-ordered relations.
\end{proposition}
\begin{comments}
  \item This can also be shown as a consequence of \cref{thm:cardinal_ordering_compatibility}, however we prefer explicit reasoning here.
\end{comments}
\begin{proof}
  \Cref{thm:set_domination_is_preorder} implies that cardinal ordering is a preorder. Antisymmetry follows from uniqueness --- if \( \kappa \) and \( \mu \) dominate each other, \fullref{thm:cantor_schroder_bernstein_theorem} implies that they are equinumerous, and \cref{thm:initial_ordinal_existence} shows that, as equinumerous initial ordinals, they are equal.

  \Fullref{thm:set_domination_is_connected} implies that cardinal ordering is connected, so it is a total order.

  It remains to show that it is well-founded. Indeed, fix a nonempty set \( A \) of cardinals. Let \( \mu \) be the smallest with respect to set membership (note that \( A \) is well-founded as a set of ordinals irrespective of the axiom of foundation). Then \( \mu \in \nu \) for every \( \nu \neq \mu \) in \( A \), thus \( \mu \subseteq \nu \) since \( \nu \) is transitive, inclusion \( \iota: \mu \to \nu \) shows that \( \nu \) dominates \( \mu \).
\end{proof}

\begin{proposition}\label{thm:cardinal_ordering_compatibility}
  The \hyperref[def:cardinal_ordering]{ordering of cardinals} coincides with the \hyperref[def:ordinal_ordering]{ordering} of their \hyperref[def:initial_ordinal]{initial ordinals}.
\end{proposition}
\begin{comments}
  \item As a consequence of \cref{thm:ordinal_ordering_is_well_ordering}, in \hyperref[def:zfc]{\logic{ZFC}}, the axioms for well-ordered relations can be derived for cardinal ordering.
\end{comments}
\begin{proof}
  Let \( \alpha \) and \( \beta \) be the initial ordinals of \( \kappa \) and \( \mu \), respectively. We must show that
  \begin{equation*}
    \alpha \in \beta \T{if and only if} \kappa \leq \mu.
  \end{equation*}

  \SufficiencySubProof If \( \alpha \in \beta \), then \( \alpha \subseteq \beta \) since \( \beta \) is transitive, so we can an inclusion function \( \iota: \alpha \to \beta \). Since \( \kappa \) and \( \alpha \) are formally identical, as are \( \mu \) and \( \beta \), this inclusion is an injective function from \( \kappa \) to \( \mu \).

  \NecessitySubProof Suppose that there exists an injective function \( f: \kappa \to \mu \). By \cref{thm:ordinal_trichotomy}, \( \alpha \) and \( \beta \) are comparable.

  If \( \beta \in \alpha \), then \cref{thm:ordinal_ordering_is_well_ordering} implies that \( \beta \subsetneq \alpha \), so there cannot be an injective function from \( \alpha \) to \( \beta \). But \( f \) is such a function, so \( \beta \) cannot be a member of \( \alpha \).

  This leaves us with either \( \alpha \in \beta \) or \( \alpha = \beta \).
\end{proof}

\begin{corollary}\label{thm:lattice_of_cardinals}
  The union of a set of \hyperref[def:cardinal]{cardinal} is their \hyperref[def:extremal_points/supremum_and_infimum]{supremum} (with respect to \hyperref[def:set_domination]{set domination}, while the intersection of a nonempty set of cardinals is their \hyperref[def:extremal_points/maximal_and_minimal_element]{minimum}.
\end{corollary}
\begin{proof}
  \Cref{thm:def:initial_ordinal/union} implies that the union of cardinals is itself a cardinal, and their supremum with respect to set membership. \Cref{thm:cardinal_ordering_compatibility} implies that it is also a supremum with respect to set domination.

  \Cref{thm:def:ordinal_ordering/minimum} implies that the intersection of cardinals is their minimum with respect to set membership, so it is also their minimum with respect to set domination.
\end{proof}

\begin{proposition}\label{thm:hartogs_number_is_successor_cardinal}
  The \hyperref[def:hartogs_number]{Hartogs' number} of a \hyperref[def:cardinal]{cardinal} is the smallest cardinal that \hyperref[def:set_domination]{strictly dominates} it.
\end{proposition}
\begin{proof}
  \Fullref{thm:hartogs_lemma} implies that the Hartogs' number \( \alpha \) of \( \kappa \) is the smallest ordinal that strictly dominates it, and \cref{thm:hartogs_number_is_initial_ordinal} implies that \( \alpha \) is in fact a cardinal.
\end{proof}

\begin{definition}\label{def:cardinal_successor}\mcite[20]{Jech2006SetTheory}
  \Cref{thm:hartogs_number_is_successor_cardinal} suggests defining the \term[en=successor (cardinal) (\cite[def. 6.2.24(i)]{Hinman2005Logic})]{successor} of a \hyperref[def:cardinal]{cardinal} as its \hyperref[def:hartogs_number]{Hartogs' number}.
\end{definition}
\begin{comments}
  \item We use the cardinal successor operation to characterize successors cardinals in \cref{def:successor_and_limit_cardinal}.
\end{comments}

\begin{proposition}\label{thm:successor_cardinal_and_power_set}
  The \hyperref[def:power_set]{power set} \( \pow(\kappa) \) of a \hyperref[def:cardinal]{cardinal} \( \kappa \) \hyperref[def:set_domination]{dominates} the \hyperref[def:successor_and_limit_cardinal]{successor} of \( \kappa \).
\end{proposition}
\begin{comments}
  \item The converse cannot be proved in \hyperref[def:zfc]{\logic{ZFC}} --- see \cref{rem:hyp:continuum_hypothesis}.
\end{comments}
\begin{proof}
  \Fullref{thm:cantors_power_set_theorem} implies that \( \pow(\kappa) \) strictly dominates \( \kappa \), so its cardinality \( \mu \) also strictly dominates \( \kappa \).

  \Cref{thm:hartogs_number_is_successor_cardinal} implies that the successor cardinal \( \sigma \) of \( \kappa \) is the smallest one that strictly dominates it. So \( \mu \geq \sigma > \kappa \).
\end{proof}

\paragraph{Finite and infinite sets}

\begin{proposition}\label{thm:members_of_omega_are_cardinals}
  The natural numbers, as members of \hyperref[def:smallest_inductive_set]{smallest inductive set} \( \omega \), are \hyperref[def:cardinal]{cardinals}.
\end{proposition}
\begin{proof}
  Restatement of \cref{thm:members_of_omega_are_initial_ordinals}.
\end{proof}

\begin{definition}\label{def:aleph_hierarchy}\mcite[30]{Jech2006SetTheory}
  We have defined the \( \omega \) hierarchy of ordinals in \cref{def:omega_hierarchy}, and we have shown in \cref{thm:omega_hierarchy_ordinals_iff_initial} that its elements are precisely the \hyperref[def:initial_ordinal]{initial ordinals} containing \( \omega_0 \) as a subset.

  We can thus regard \( \omega_\alpha \) as a \hyperref[def:cardinal]{cardinal}, in which case we use the notation \( \aleph_\alpha \) instead of \( \omega_\alpha \). We can restate the definition as follows:
  \begin{equation*}
    \aleph_\alpha \coloneqq \begin{cases}
      \omega_0,                                                         &\alpha = \varnothing, \\
      \T{\hyperref[def:cardinal_successor]{successor} of} \aleph_\beta, &\alpha = S(\beta), \\
      \bigcup\set{ \aleph_\beta \given \beta \in \alpha }               &\alpha \T{is a nonzero limit ordinal.}
    \end{cases}
  \end{equation*}
\end{definition}
\begin{comments}
  \item This hierarchy is important because it describes all infinite cardinals, as shown in \cref{thm:cardinal_iff_finite_or_aleph}. It is connected to the simpler \hyperref[def:beth_hierarchy]{\( \beth \) hierarchy} via \fullref{hyp:generalized_continuum_hypothesis}.
\end{comments}

\begin{proposition}\label{thm:def:aleph_hierarchy}
  The \( \aleph \) hierarchy from \cref{def:aleph_hierarchy} has the following basic properties:
  \begin{thmenum}
    \thmitem{thm:def:aleph_hierarchy/monotone} It is monotone: \( \beta \in \alpha \) implies that \( \aleph_\beta < \aleph_\alpha \).
  \end{thmenum}
\end{proposition}
\begin{proof}
  \SubProofOf{thm:def:aleph_hierarchy/monotone} Restatement of \cref{thm:def:omega_hierarchy/monotone}.
\end{proof}

\begin{definition}\label{def:aleph_zero_related_cardinalities}
  In \cref{tab:def:aleph_zero_related_cardinalities}, we define terminology for \hyperref[def:cardinality]{cardinalities} with respect to the first cardinal \( \aleph_0 \) from \cref{def:aleph_hierarchy}.

  \begin{table}
    \begin{center}
      \begin{tabular}{ll}
        \toprule
        Cardinality         & Terminology \\
        \midrule
        \( < \aleph_0 \)    & \term[ru=конечное (множество) (\cite[154]{КолмогоровДрагалин2006ЛогикаДополнительныеГлавы}), en=finite (set) (\cite[def. 6.2.32(iii)]{Hinman2005Logic})]{Finite} \\
        \( \leq \aleph_0 \) & \term[en=at most countable (set) (\cite[30]{Jech2006SetTheory})]{At most countable} \\
        \( = \aleph_0 \)    & \term[en=countably infinite (set) (\cite[def. I.10.5]{Kunen2013SetTheory})]{Countably infinite} \\
        \( > \aleph_0 \)    & \term[bg=неизброимо (множество) (\cite[225]{ГеновМиховскиМоллов1991Алгебра}), ru=несчётное (множество) (\cite[14]{Александров1977ОбщаяТопология}), en=uncountable (set) (\cite[def. I.10.5]{Kunen2013SetTheory})]{Uncountable} \\
        \( \geq \aleph_0 \) & \term[ru=бесконечное (множество) (\cite[154]{КолмогоровДрагалин2006ЛогикаДополнительныеГлавы}), en=infinite (set) (\cite[def. 6.2.32(iii)]{Hinman2005Logic})]{Infinite} \\
        \bottomrule
      \end{tabular}
    \end{center}
    \caption{Terminology based on set cardinality.}\label{tab:def:aleph_zero_related_cardinalities}
  \end{table}
\end{definition}
\begin{comments}
  \item See \cref{rem:countability_terminology} some a discussion of ambiguity regarding some of these terms.
\end{comments}

\begin{remark}\label{rem:countability_terminology}
  We mentioned in \cref{def:aleph_zero_related_cardinalities} that the terms based on the word \enquote{countable} and \enquote{enumerable} are inconsistently used across the literature. We list in \cref{tab:rem:countability_terminology} how this is handled in some of the books we cite.

  \begin{table}
    \centering
    \begin{tabular}{l l l}
      \toprule
      Source                                                          & \( \leq \aleph_0 \)        & \( = \aleph_0 \)                \\
      \midrule
      \cite[30]{Jech2006SetTheory}                                    & At most countable          & Countable                       \\
      \cite[def. 2.3]{Rudin1976AnalysisPrinciples}                    & At most countable          & Countable                       \\
      \cite[def. I.10.5]{Kunen2013SetTheory}                          & Countable                  & Countably infinite              \\
      \cite[3]{Engelking1989GeneralTopology}                          & Countable                  & Countably infinite              \\
      \cite[25]{Kelley1975GeneralTopology}                            & Countable                  & Countably infinite              \\
      \cite[def. 6.2.32]{Hinman2005Logic}                             & Countable                  & Denumerable                     \\
      \cite[xxii]{Mendelson2015Logic}                                 & Countable                  & Denumerable                     \\
      \cite[3]{Harzheim2005OrderedSets}                               & Countable                  & Denumerable                     \\
      \cite[91]{Halmos1960NaiveSetTheory}                             & Countable / denumerable    & Countably infinite              \\
      \cite[176]{Kleene2002Logic}                                     & Countable / enumerable     & Countably / enumerably infinite \\
      \cite[159]{Enderton1977SetTheory}                               & Countable                  & ---                             \\
      \cite[13]{ШеньВерещагин2020НачалаТеорииМножеств}                & ---                        & Счётное (Countable)             \\
      \cite[def. 1.3]{Александров1977ОбщаяТопология}                  & ---                        & Счётное (Countable)             \\
      \bottomrule
    \end{tabular}
    \caption{The table for \cref{rem:countability_terminology} discussing terminology related to \hyperref[def:aleph_zero_related_cardinalities]{countability}.}\label{tab:rem:countability_terminology}
  \end{table}
\end{remark}

\begin{remark}\label{rem:cardinality_auxiliary_terminology}
  There are multiple occasions where we have to use adjectives based on cardinality. The simplest example is when we want to call a set \enquote{finite} (in accordance with \cref{def:aleph_zero_related_cardinalities}), \enquote{countably infinite} (in accordance with \cref{def:aleph_zero_related_cardinalities}), \enquote{uncountable}, etc.

  As another example, a \hyperref[rem:graph_of_unspecified_type]{graph or hypergraph} can be \enquote{locally finite}, \enquote{locally countable}, etc., in accordance with \cref{def:graph_cardinality/local}.

  On the other hand, the graph itself can have finitely many arcs/edges but uncountably many vertices. So we call a graph \enquote{finite} if both the vertices and edges are finite. So even though we use the terminology for cardinals, we are not dealing with the cardinality of a single set. As a general rule, if some contraption is defined as an \hyperref[def:ordered_tuple]{ordered tuple} \( \mscrA = (A_1, \ldots, A_n) \) of sets, when using cardinality-related terminology for \( \mscrA \), we refer to the largest among \( \card(A_k) \).

  In order not to repeat the same list of adjectives in every relevant definition, we reuse the adjectives from \cref{tab:def:aleph_zero_related_cardinalities}.
\end{remark}
\begin{comments}
  \item Note that even if a set is finite, it may not be hereditarily finite, i.e. it may consist of infinite sets. This somewhat contrasts with the concepts of \enquote{small} and \enquote{large} sets from \cref{def:small_set}, where, due to \cref{thm:small_set_closure}, all elements of a small set are small, and all components of a small ordered tuple are small.
\end{comments}

\begin{proposition}\label{thm:cardinal_is_finite_iff_natural_number}
  A \hyperref[def:cardinal]{cardinal} is \hyperref[def:aleph_zero_related_cardinalities]{finite} if and only if it is a \hyperref[def:natural_numbers]{natural number} (i.e. a member of \( \aleph_0 \)).
\end{proposition}
\begin{proof}
  \Cref{thm:cardinal_ordering_compatibility} implies that if \( \kappa \) is finite, i.e. if \( \kappa < \aleph_0 \) with respect to set domination, then \( \kappa \in \aleph_0 \).

  Conversely, \cref{thm:members_of_omega_are_cardinals} implies that all members of \( \aleph_0 \) are cardinals, so \( \kappa \in \aleph_0 \) entails \( \kappa < \aleph_0 \).
\end{proof}

\begin{proposition}\label{thm:cardinal_is_infinite_iff_aleph}
  A \hyperref[def:cardinal]{cardinal} is \hyperref[def:aleph_zero_related_cardinalities]{infinite} if and only if it belongs to the \( \aleph \) hierarchy from \cref{def:aleph_hierarchy}.
\end{proposition}
\begin{proof}
  If a cardinal \( \kappa \) is infinite, i.e. if \( \kappa \geq \aleph_0 \), \cref{thm:omega_hierarchy_ordinals_iff_initial} implies that, as an initial ordinal, it belongs to the \( \omega \) hierarchy (which we have used to define the \( \aleph \) hierarchy).

  Conversely, every \( \aleph_\alpha \) is a cardinal by \cref{thm:omega_hierarchy_ordinals_iff_initial} and, furthermore, \( \aleph_0 \subseteq \aleph_\alpha \) because \cref{thm:def:omega_hierarchy/monotone} implies that \( \aleph_0 \in \aleph_\alpha \) whenever \( \varnothing \in \alpha \).
\end{proof}

\begin{corollary}\label{thm:cardinal_iff_finite_or_aleph}
  Every \hyperref[def:cardinal]{cardinal} is either a \hyperref[def:natural_numbers]{natural number} or a member of the \( \aleph \) hierarchy from \cref{def:aleph_hierarchy}.
\end{corollary}
\begin{proof}
  The finite case follows from \cref{thm:cardinal_is_finite_iff_natural_number} and the infinite case follows from \cref{thm:cardinal_is_infinite_iff_aleph}.
\end{proof}

\begin{definition}\label{def:cardinality_of_the_continuum}\mcite[37]{Jech2006SetTheory}
  We denote the \hyperref[def:cardinality]{cardinality} of the \hyperref[def:power_set]{power set} \( \pow(\aleph_0) \) by \( \mfrakc \) and call it the \term[ru=мощность континуума (\cite[33]{Александров1977ОбщаяТопология})]{cardinality of the continuum}.
\end{definition}
\begin{comments}
  \item Using cardinal exponentiation, which we will define in \cref{def:cardinal_arithmetic/exponentiation}, we can express this as \( \mfrakc = 2^{\aleph_0} \).
\end{comments}
\begin{comments}
  \item We prove in \cref{thm:cardinality_of_the_continuum} that the cardinality of the \hyperref[def:real_numbers]{real numbers} is \( c \).
\end{comments}

\begin{conjecture}[Continuum hypothesis]\label{hyp:continuum_hypothesis}\mcite[165]{Enderton1977SetTheory}
  The \hyperref[def:cardinality_of_the_continuum]{cardinality of the continuum} \( c \) equals the \hyperref[def:aleph_zero_related_cardinalities]{first uncountable cardinal} \( \aleph_1 \).
\end{conjecture}
\begin{comments}
  \item \Cref{thm:successor_cardinal_and_power_set} implies \( \aleph_1 \leq c \), however the converse is not provable. The hypothesis and its stronger form \fullref{hyp:generalized_continuum_hypothesis} are independent of \hyperref[def:zfc]{\logic{ZFC}} --- see \cref{rem:hyp:continuum_hypothesis}.
\end{comments}

\begin{remark}\label{rem:hyp:continuum_hypothesis}
  \Textcite[166]{Enderton1977SetTheory} summarizes the history of the continuum hypothesis as follows\fnote{By \enquote{our axioms of set theory}, Enderton means \hyperref[def:zfc]{\( \logic{ZFC} \)}.}:
  \begin{displayquote}
    Cantor conjectured that the continuum hypothesis was true. And David Hilbert later published a purported proof. But the proof was incorrect, and more recent work has cast doubt on the continuum hypothesis. In 1939 G\"odel proved that on the basis of our axioms for set theory (which we here assume to be consistent) the continuum hypothesis could not be disproved. Then in 1963 Paul Cohen showed that the continuum hypothesis could not be proved from our axioms either.
  \end{displayquote}
\end{remark}

\begin{theorem}[Incompleteness of ZFC]\label{thm:zfc_incompleteness}
  \hyperref[def:zfc]{Zermelo-Fraenkel set theory} is \hyperref[def:fol_theory/complete]{incomplete}.
\end{theorem}
\begin{proof}
  As discussed in \cref{rem:hyp:continuum_hypothesis}, \fullref{hyp:continuum_hypothesis} can neither be proved nor disproved in \logic{ZFC}.
\end{proof}

\begin{conjecture}[Generalized continuum hypothesis]\label{hyp:generalized_continuum_hypothesis}\mcite[166]{Enderton1977SetTheory}
  The \hyperref[def:ordinal_successor]{successor} of every \hyperref[def:cardinal]{cardinal} \( \kappa \) is \hyperref[def:cardinality]{cardinality} of its \hyperref[def:power_set]{power set} \( \pow(\kappa) \).
\end{conjecture}
\begin{comments}
  \item This is in general only an equality --- see \cref{thm:aleph_and_beth_hierarchy}.
\end{comments}

\begin{definition}\label{def:beth_hierarchy}\mcite[214]{Enderton1977SetTheory}
  Similarly to the \( \aleph \) hierarchy from \cref{def:aleph_hierarchy}, we can define the following hierarchy:
  \begin{equation*}
    \beth_\alpha \coloneqq \begin{cases}
      \omega_0,                                          &\alpha = \varnothing \\
      \card(\pow(\beta)),                                &\alpha = S(\beta) \\
      \bigcup\set{ \beth_\beta \given \beta \in \alpha } &\alpha \T{is a nonzero limit ordinal.}
    \end{cases}
  \end{equation*}
\end{definition}

\begin{proposition}\label{thm:aleph_and_beth_hierarchy}
  In the \( \aleph \) and \( \beta \) hierarchies, for every index \( \alpha \) we have
  \begin{equation*}
    \aleph_\alpha \leq \beta_\alpha.
  \end{equation*}
\end{proposition}
\begin{comments}
  \item The converse is only true assuming \fullref{hyp:generalized_continuum_hypothesis}.
\end{comments}
\begin{proof}
  Follows from \cref{thm:successor_cardinal_and_power_set}.
\end{proof}

\begin{definition}\label{def:dedekind_finiteness}\mcite[\S 6.27]{Schechter1997AnalysisHandbook}
  We say that a set \term{Dedekind infinite} if it is \hyperref[def:equinumerosity]{equinumerous} with a proper subset of itself. Otherwise, we say that it is \term{Dedekind finite}.
\end{definition}

\begin{lemma}\label{thm:cardinal_infinite_iff_dedekind_infinite}
  A \hyperref[def:cardinal]{cardinal} is \hyperref[def:aleph_zero_related_cardinalities]{infinite} if and only if it is \hyperref[def:dedekind_finiteness]{Dedekind infinite}.
\end{lemma}
\begin{proof}
  \SufficiencySubProof Suppose that \( \kappa \) is an infinite cardinal \( \kappa \), i.e. \( \kappa \geq \aleph_0 \). Let
  \begin{equation*}
    \begin{aligned}
      &f: \kappa \setminus \set{ \varnothing } \to \kappa, \\
      &f(\alpha) \coloneqq \begin{cases}
        \beta,  &\alpha = S(\beta) \in \omega, \\
        \alpha, &\T{otherwise.} \\
      \end{cases}
    \end{aligned}
  \end{equation*}

  We claim that \( f \) is bijective.

  \SubProofOf*[def:function_invertibility/injective]{injectivity} If \( f(\alpha) = f(\alpha') \), we have several possibilities:
  \begin{itemize}
    \item If \( \alpha \in \omega \) and \( \alpha' \in \omega \), then \( \alpha = S(\beta) \) and \( \alpha' = S(\beta') \) and
    \begin{equation*}
      \beta = f(\alpha) = f(\alpha') = \beta',
    \end{equation*}
    so
    \begin{equation*}
      \alpha = S(\beta) = S(\beta') = \alpha'.
    \end{equation*}

    \item If \( \alpha \in \omega \) and \( \alpha' \not\in \omega \), then \( \alpha = S(\beta) \) and
    \begin{equation*}
      \beta = f(\alpha) = f(\alpha') = \alpha',
    \end{equation*}
    which implies that \( \beta \not\in \omega \) --- a contradiction.

    \item If \( \alpha \not\in \omega \) and \( \alpha' \in \omega \), we similarly reach a contradiction.

    \item If \( \alpha \not\in \omega \) and \( \alpha' \not\in \omega \), then
    \begin{equation*}
      \alpha = f(\alpha) = f(\alpha') = \alpha'.
    \end{equation*}
  \end{itemize}

  \SubProofOf*[def:function_invertibility/surjective]{surjectivity} For any \( \alpha \) in \( \kappa \), we have two possibilities:
  \begin{itemize}
    \item If \( \alpha \in \omega \), then \( f(S(\alpha)) = \alpha \).
    \item If \( \alpha \not\in \omega \), then \( f(\alpha) = \alpha \).
  \end{itemize}

  \NecessitySubProof Suppose that \( \kappa \) is a Dedekind infinite cardinal. \Cref{thm:natural_number_is_not_equinumerous_to_proper_subset} implies that the natural numbers are Dedekind finite, so \( \kappa \) cannot be a natural number. It remains for \( \kappa \) to be infinite.
\end{proof}

\begin{proposition}\label{thm:infinite_iff_dedekind_infinite}
  A set is \hyperref[def:aleph_zero_related_cardinalities]{infinite} if and only if it is \hyperref[def:dedekind_finiteness]{Dedekind infinite}.
\end{proposition}
\begin{proof}
  Fix a set \( A \). If \( \kappa \) is its cardinality, there exists a corresponding bijective function \( f: A \to \kappa \).

  \SufficiencySubProof Suppose that \( A \) is infinite, i.e. \( \kappa \geq \aleph_0 \). \Cref{thm:cardinal_infinite_iff_dedekind_infinite} implies that \( \kappa \) is Dedekind infinite, so it has a proper subset \( K \) for which there exists a bijective function \( g: \kappa \to K \).

  Then \( B \coloneqq f^{-1}[K] \) is a strict subset of \( A \), and
  \begin{equation*}
    f^{-1}\restr_{K,B} \bincirc g \bincirc f: A \to B
  \end{equation*}
  is bijective.

  Thus, \( A \) is Dedekind infinite.

  \NecessitySubProof Conversely, if \( A \) is Dedekind infinite, there exists a bijective function \( g: A \to B \) into a proper subset \( B \) of \( A \).

  Then \( K \coloneqq f[B] \) is a strict subset of \( \kappa \), and
  \begin{equation*}
    f\restr_{B,K} \bincirc g \bincirc f^{-1}: \kappa \to K
  \end{equation*}
  is bijective.

  Thus, \( \kappa \) is Dedekind infinite, and \cref{thm:cardinal_infinite_iff_dedekind_infinite} implies that it is infinite, i.e. \( \kappa \geq \aleph_0 \).
\end{proof}

\paragraph{Successor and limit cardinals}

\begin{definition}\label{def:successor_and_limit_cardinal}
  We say that a \hyperref[def:cardinal]{cardinal} is a \term{successor cardinal} if it is the \hyperref[def:cardinal_successor]{successor} (i.e. the \hyperref[def:hartogs_number]{Hartogs' number}) of another cardinal. Otherwise, we say that it is a \term{limit cardinal}.
\end{definition}
\begin{comments}
  \item These notions are analogous to successor and limit ordinals, defined in \cref{def:successor_and_limit_ordinal}. Even though we have defined cardinals in \cref{def:cardinal} as initial ordinals, \cref{thm:cardinal_is_infinite_iff_limit_ordinal} implies that infinite successor cardinals are limit ordinals, so the two notions are fundamentally distinct.

  \item See \cref{thm:limit_cardinal_equivalences} for some equivalent conditions.

  \item An important special case of limit cardinals are \enquote{strong} limits cardinals --- see \cref{def:strong_limit_cardinal}.
\end{comments}

\begin{lemma}\label{thm:successor_ordinal}
  If a \hyperref[def:successor_and_limit_ordinal]{successor ordinal} \( \alpha = S(\beta) \) is \hyperref[def:dedekind_finiteness]{Dedekind infinite}, then \( \alpha \) and \( \beta \) are \hyperref[def:equinumerosity]{equinumerous}.
\end{lemma}
\begin{proof}
  Suppose that \( \alpha = S(\beta) \) has a proper subset \( A \), for which there exists a bijective function \( f: A \to \alpha \).

  If \( \beta \) belongs to \( A \), we can instead fix some \( \beta' \) from \( \alpha \setminus A \) and define \( A' \coloneqq (A \setminus \set{ \beta }) \cup \set{ \beta' } \) and
  \begin{equation*}
    \begin{aligned}
      &f': A' \to \alpha, \\
      &f'(\gamma) \coloneqq \begin{cases}
        f(\beta),  &\gamma = \beta', \\
        f(\gamma), &\T{otherwise.}
      \end{cases}
    \end{aligned}
  \end{equation*}

  Otherwise, let \( \beta' \), \( A' \) and \( f' \) match \( \beta \), \( A \) and \( f \).

  Since \( \alpha = S(\beta) = \beta \cup \set{ \beta } \) and since \( \beta \) is not in \( A' \), we conclude that \( A' \) is a subset of \( \beta \). Therefore, \( \beta \) dominates \( A' \), which is equinumerous with \( \alpha \). But \( \beta \) is also a subset of \( \alpha \), so \( \alpha \) dominates \( \beta \). By \fullref{thm:cantor_schroder_bernstein_theorem}, \( \alpha \) and \( \beta \) are equinumerous.
\end{proof}

\begin{proposition}\label{thm:cardinal_is_finite_iff_successor_ordinal}
  A nonzero \hyperref[def:cardinal]{cardinal} is \hyperref[def:aleph_zero_related_cardinalities]{finite} if and only if it is a \hyperref[def:successor_and_limit_ordinal]{successor ordinal}.
\end{proposition}
\begin{proof}
  \SufficiencySubProof If \( 0 < \kappa < \alpha_0 \), \cref{thm:cardinal_ordering_compatibility} implies that \( \varnothing \in \kappa \in \omega \), and \cref{thm:omega_is_limit_ordinal} implies that \( \kappa \) is a successor ordinal.

  \NecessitySubProof Suppose that \( \kappa = S(\beta) \) is a successor ordinal. Suppose that \( \kappa \) is \hyperref[def:aleph_zero_related_cardinalities]{infinite}. Then \cref{thm:cardinal_infinite_iff_dedekind_infinite} implies that \( \kappa \) is \hyperref[def:dedekind_finiteness]{Dedekind infinite}, and \cref{thm:successor_ordinal} implies that \( \kappa \) and \( \beta \) are \hyperref[def:equinumerosity]{equinumerous}. But \( \kappa \) is a cardinal, hence an initial ordinal, so it cannot be equinumerous with its member \( \beta \).

  The obtained contradiction shows that \( \kappa \) cannot be infinite.
\end{proof}

\begin{corollary}\label{thm:cardinal_is_infinite_iff_limit_ordinal}
  A nonzero \hyperref[def:cardinal]{cardinal} is \hyperref[def:aleph_zero_related_cardinalities]{infinite} if and only if it is a \hyperref[def:successor_and_limit_ordinal]{limit ordinal}.
\end{corollary}
\begin{proof}
  Dual of \cref{thm:cardinal_is_finite_iff_successor_ordinal}.
\end{proof}

\begin{proposition}\label{thm:limit_cardinal_equivalences}
  The cardinal \( \aleph_\alpha \) is a \hyperref[def:successor_and_limit_cardinal]{limitz cardinal} if and only if \( \alpha \) is a \hyperref[def:successor_and_limit_ordinal]{limit ordinal}.
\end{proposition}
\begin{proof}
  Trivial.
\end{proof}

\paragraph{Closure under operations}

\begin{proposition}\label{thm:finite_unions_and_products_are_finite}
  All \hyperref[def:aleph_zero_related_cardinalities]{finite} \hyperref[def:set_union]{unions} and \hyperref[def:cartesian_product]{Cartesian products} of finite sets are finite.
\end{proposition}
\begin{proof}
  Let \( \mscrA \) be a finite set of finite sets. Let \( n \) be the cardinality of \( \mscrA \) and let \( m_A \) be the cardinality of \( A \in \mscrA \). \Cref{thm:def:extremal_points/finite_set_has_maximal_element} implies that there is a maximal cardinality \( m \) among these \( m_A \).

  Then there exist some bijective functions \( f: \mscrA \to n \) and \( g_A: A \to m_A \) for every \( A \) in \( \mscrA \).

  \SubProof{Proof for Cartesian products} Consider the function
  \begin{equation*}
    \begin{aligned}
      &h: \prod_{A \in \mscrA} A \to n \times m, \\
      &h(\seq{ a_A }_{A \in \mscrA}) \coloneqq (f(A), g_A(a)).
    \end{aligned}
  \end{equation*}

  It is clearly injective.

  \SubProof{Proof for unions} Let \( c: \bigcup \mscrA \to \mscrA \) be a function picking, for every \( a \) in the union, a set \( A \) containing \( a \). For the sake of determinism, we can choose \( A \) that minimizes \( f(A) \). Then we can define the following:
  \begin{equation*}
    \begin{aligned}
      &h: \bigcup \mscrA \to n \times m \\
      &h(a) \coloneqq (f(c_a), g_{c_a}(a)),
    \end{aligned}
  \end{equation*}

  It is also clearly injective.
\end{proof}

\begin{proposition}\label{thm:omega_equinumerous_with_omega_squared}
  The smallest inductive set \( \omega \) is equinumerous with \( \omega \times \omega \).
\end{proposition}
\begin{proof}
  We can give a short proof using \fullref{thm:cantor_schroder_bernstein_theorem} using the injective functions \( (n, m) \mapsto 2^n 3^m \) in one direction and \( n \mapsto (n, 0) \) in the other direction. Proving the injectivity of \( f \), however, requires \fullref{thm:fundamental_theorem_of_arithmetic}, and we prove the latter using machinery from \fullref{ch:ring_theory}. We will instead give a direct proof with an explicit construction. We will construct a bijective function from \( \omega \times \omega \) to \( \omega \) --- the function visualized in \cref{fig:thm:omega_equinumerous_with_omega_squared}.

  We begin by defining the diagonal in \cref{fig:thm:omega_equinumerous_with_omega_squared}. For each natural number \( k \), define the set of pairs that sum to \( k \):
  \begin{equation*}
    A_k \coloneqq \set{ (n, m) \in \omega \times \omega \given n + m = k }.
  \end{equation*}

  We can use induction to show that \( \card(A_k) = k + 1 \). That is, \( A_k \) can \enquote{fit} \( k + 1 \) numbers. We can now define the function
  \begin{equation*}
    \begin{aligned}
      &d: \omega \to \omega \\
      &d(n) \coloneqq \sum_{k=0}^n \card(A_n)
    \end{aligned}
  \end{equation*}
  that gives us how many numbers we have already \enquote{fit} in the first \( n \) diagonals.

  It is clear that the point \( (n, m) \) lies in \( A_{n + m} \). We want to know how many numbers we have \enquote{fit} in the diagonal prior to that, for which we can use \( d(n + m - 1) \). This leads us to the definition
  \begin{equation*}
    \begin{aligned}
      &f: \omega \times \omega \to \omega \\
      &f(n, m) \coloneqq \begin{cases}
        0,                &n + m = 0 \\
        d(n + m - 1) + n, &n + m > 0. \\
      \end{cases}
    \end{aligned}
  \end{equation*}

  We will first show that \( f \) is injective using induction on \( n + m \) (that is, on the diagonals). Suppose that \( f(n_1, m_1) = f(n_2, m_2) \) implies \( n_1 = n_2 \) and \( m_1 = m_2 \) for all pairs with a sum less than \( l \). Let \( (n_1, m_1) \) and \( (n_2, m_2) \) be two points in \( A_l \) such that \( f(n_1, m_1) = f(n_2, m_2) \). The cases \( l <= 1 \) are trivial, so suppose that \( l > 1 \). Then
  \begin{equation*}
    f(n_1, m_1)
    =
    d(n_1 + m_1 - 1) + n_1
    =
    l - 1 + d(n_1 + m_1 - 2) + n_1
    =
    l + f(n_1 - 1, m_1).
  \end{equation*}

  We can now apply the inductive hypothesis and obtain that \( n_1 = n_2 \) and \( m_1 = m_2 \). This proves injectivity.

  To see that \( f \) is surjective, we will use induction on \( k \in \omega \). The base case is again trivial. Now suppose that \( n + m > 0 \) and
  \begin{equation*}
    f(n, m) = d(n + m - 1) + n = k.
  \end{equation*}

  We have two cases:
  \begin{itemize}
    \item If \( m = 0 \), then
    \begin{equation*}
      f(0, n + m + 1) = d(n + m) + 0 = d(n + m - 1) + n + m = f(n, m) + 1.
    \end{equation*}

    \item If \( m > 0 \), then
    \begin{equation*}
      f(n + 1, m - 1) = d(n + m - 1) + n + 1 = f(n, m) + 1.
    \end{equation*}
  \end{itemize}

  In both cases we have shown that \( k + 1 = f(n, m) + 1 \) is in the image of \( f \), which concludes our proof of surjectivity (and hence bijectivity).

  Lastly, although it is not necessary for the proof, we can expand the definition of \( d \) to see that \( f \) is actually a \hyperref[def:polynomial_algebra/polynomials]{polynomial}:
  \begin{equation*}
    f(n, m)
    =
    \sum_{k=0}^{n + m - 1} (k + 1) + n
    =
    \sum_{k=1}^{n + m} k + n
    \reloset {\eqref{eq:thm:numeric_arithmetic_progression_partial_sums}} =
    \frac {(n + m) (n + m + 1)} 2 + n.
  \end{equation*}

  As an added benefit, this polynomial also handles the case \( n = m = 0 \).

  \begin{figure}
    \hfill
    \includegraphics[page=1]{output/thm__omega_equinumerous_with_omega_squared}
    \hfill
    \includegraphics[page=2]{output/thm__omega_equinumerous_with_omega_squared}
    \hfill\hfill
    \caption{Visualization on an integer coordinate grid of the diagonal sets \( A_k \) and of the bijective function defined in \cref{thm:omega_equinumerous_with_omega_squared}.}\label{fig:thm:omega_equinumerous_with_omega_squared}
  \end{figure}
\end{proof}

\begin{corollary}\label{thm:countable_product_of_countable_sets}
  A finite \hyperref[def:cartesian_product]{Cartesian product} of at most countable sets is at most countable.
\end{corollary}
\begin{proof}
  Let \( A_1, \ldots, A_n \) be a finite family of at most countable sets.

  Suppose that the inductive hypothesis holds for \( n \). Countability ensures that there exist injective functions \( g: A_1 \times \cdots \times A_n \times \omega \) and \( h: A_{n+1} \to \omega \). Denote by \( f \) the bijective function from \( \omega \) to \( \omega \times \omega \) obtained in \cref{thm:omega_equinumerous_with_omega_squared} and define
  \begin{equation*}
    \begin{aligned}
      &F: A_1 \times \cdots \times A_n \times A_{n+1} \to \omega \\
      &F(a_1, \ldots, a_n, a_{n+1}) \coloneqq f(g(a_1, \ldots, a_n), h(a_{n+1}))
    \end{aligned}
  \end{equation*}

  By \cref{thm:function_superposition_invertibility}, the function \( F \) is injective as a \hyperref[def:function_superposition]{superposition} of injective functions.

  Therefore, \( \omega \) dominates the product \( A_1 \times \cdots \times A_n \), i.e. the product is countable.
\end{proof}

\begin{proposition}\label{thm:countably_infinite_disjoin_union_of_nonempty_sets}
  A \hyperref[def:aleph_zero_related_cardinalities]{countably infinite} \hyperref[def:disjoint_union]{disjoint union} of nonempty sets is infinite.
\end{proposition}
\begin{comments}
  \item \todo{Consider stating as a special case of \cref{thm:cardinality_of_disjoint_union}}.
\end{comments}
\begin{proof}
  Let \( \seq{ A_k }_{k \in \omega} \) be a countably infinite family of nonempty sets and consider the family \( \seq{ B_k }_{k \in \omega} \), where
  \begin{equation*}
    B_k \coloneqq \set{ \iota_k(a) \given a \in A_k }.
  \end{equation*}

   We can use the \hyperref[def:zfc/choice]{axiom of choice} to construct a sequence
   \begin{equation*}
     \iota_0(a_0), \iota_1(a_1), \ldots
   \end{equation*}
   of elements of the product \( \prod_{k \in \omega} B_k \).

   This is in fact an injective function from \( \omega \) to the disjoint union
   \begin{equation*}
     \bigcup_{k \in \omega} B_k = \coprod_{k \in \omega} A_k.
   \end{equation*}

   Therefore, the union is at least countably infinite.
\end{proof}

\begin{proposition}\label{thm:countably_infinite_union_of_countably_infinite_sets}\mcite[thm. 6Q]{Enderton1977SetTheory}
  A \hyperref[def:aleph_zero_related_cardinalities]{countably infinite} union of countably infinite sets is countably infinite.
\end{proposition}
\begin{comments}
  \item \todo{Consider stating as a special case of \cref{thm:cardinality_of_disjoint_union}}.
\end{comments}
\begin{proof}
  Let \( \seq{ A_k }_{k \in \omega} \) be a countably infinite family of countably infinite sets. Define instead the disjoint family
  \begin{equation*}
    B_k \coloneqq \set{ (k, a) \given a \in A_k }.
  \end{equation*}

  Denote the union of the former family by \( A \) and of the latter family by \( B \). Since each \( A_k \) is countably infinite, so is \( A \). Furthermore, there exists an obvious injective function from \( A \) to \( B \), thus
  \begin{equation}\label{eq:thm:countably_infinite_union_of_countably_infinite_sets/union_card_inequality}
    \aleph_0 \leq \card(A) \leq \card(B).
  \end{equation}

  Define the set-valued mapping
  \begin{equation*}
    \begin{aligned}
      &G: \omega \to \pow(\fun(\omega, B)) \\
      &G(k) \coloneqq \set{ g: \omega \to B_k \given g \T{is bijective} }.
    \end{aligned}
  \end{equation*}

  The value \( G(k) \) is nonempty because we have assumed that \( B_k \) is countable for every \( k \in \mscrK \). \Fullref{thm:existence_of_single_valued_selections} gives us a single-valued function \( G: \omega \to \fun(\omega, B) \). Since the family \( \seq{ B_k }_{k \in \omega} \) is disjoint, \( G \) is injective. We can thus \hyperref[con:currying]{uncurry} \( G \) to obtain a function \( g \) from \( \omega \times \omega \) to \( B \).

  To prove that \( g \) is injective, suppose that \( g(n_1, m_1) = g(n_2, m_2) \). Then
  \begin{equation*}
    G(n_1)(m_1) = g(n_1, m_1) = g(n_2, m_2) = G(n_2)(m_2).
  \end{equation*}

  Note that \( n_1 \neq n_2 \) would lead to a contradiction because  \( \seq{ B_k }_{k \in \omega} \) is a disjoint family. So \( n_1 = n_2 \) and we obtain \( m_1 = m_2 \) since the function \( G(n_1): \omega \to B_k \) is injective. Therefore, \( g \) itself is also injective.

  It is also surjective because for every \( a \in B \) there exists some \( k \in \omega \) such that
  \begin{equation*}
    a \in B_k = \img(G(k)).
  \end{equation*}

  Denote by \( f \) the bijective function from \( \omega \) to \( \omega \times \omega \) obtained in \cref{thm:omega_equinumerous_with_omega_squared}. Then the function \( f \bincirc g: \omega \to B \) is bijective by \cref{thm:function_composition_invertibility}. Therefore, the union \( B \) is countably infinite. From \eqref{eq:thm:countably_infinite_union_of_countably_infinite_sets/union_card_inequality} it follows that \( A \) is also countably infinite.
\end{proof}

\begin{corollary}\label{thm:at_most_countable_union_of_at_most_countable_sets}
  An \hyperref[def:aleph_zero_related_cardinalities]{at most countable} union of at most countable sets is at most countable.
\end{corollary}
\begin{comments}
  \item \todo{Consider stating as a special case of \cref{thm:cardinality_of_disjoint_union}}.
\end{comments}
\begin{proof}
  Let \( \seq{ A_k }_{k \in \mscrK} \) be an at most countable family of at most countable sets. Denote their union by \( A \). For every \( A_k \) let \( g_k: A_k \to \omega \) be an injective function and define the \hyperref[def:disjoint_union]{disjoint union}
  \begin{equation*}
    B_k \coloneqq A_k \amalg \set{ n \in \omega \given n \notin \img(g_k) }
  \end{equation*}
  and the bijective function
  \begin{equation*}
    \begin{aligned}
      &h_k: B_k \to \omega \\
      &h_k(x) \coloneqq \begin{cases}
        (0, g_k(x)), &x \in A_k \\
        (1, x),      &\T{otherwise.}
      \end{cases}
    \end{aligned}
  \end{equation*}

  For every \( k \in \omega \setminus \mscrK \) instead define
  \begin{equation*}
    B_k \coloneqq \set{ (k, n) \given n \in \omega }
  \end{equation*}
  and let \( h_k: B_k \to \omega \) be the obvious bijective function.

  Then
  \begin{equation*}
    A
    =
    \bigcup_{k \in \mscrK} A_k
    \subseteq
    \bigcup_{k \in \mscrK} B_k
    \subseteq
    \bigcup_{k \in \omega} B_k
  \end{equation*}
  and the latter is countably infinite by \cref{thm:countably_infinite_union_of_countably_infinite_sets}. Therefore, \( A \) is at most countable.
\end{proof}
